% --- 한 걸음씩 애니메이션 (영상 쪽: 제목 + 썸네일 + 링크만)
\begin{frame}{한 걸음씩: 클릭 $\to$ 선 이동}
\vfill
\centering
\href{run:figs/sv_train_click.mp4}{\includegraphics[height=0.72\textheight]{sv_train_click_thumb.png}}\\[0.3em]
\href{run:figs/sv_train_click.mp4}{\textcolor{ksagolddk}{▶ figs/sv\_train\_click.mp4}}
\vfill
\note{학습 한 걸음을 네 단계로 나눠 보여 주는 영상입니다. 첫째, 점마다 값을 잽니다. 점 옆 숫자가 경계까지 거리입니다. 둘째, 값이 가장 작은 점을 클릭합니다. 그 값이 지금의 m입니다. 셋째, 그 점이 밉니다. 빨간 점이면 w에 학습률 곱하기 자기 좌표를 더하고 b에 b 학습률을 더해서 선이 그 점에서 멀어집니다. 넷째, w의 길이가 1을 넘으면 w와 b를 길이로 나눕니다. 이때 선은 그대로이고 숫자만 바뀝니다. 1걸음은 20개가 모두 같은 값 0이라 다 같이 밀고, 2걸음에는 빨간 점 3.5, 6이 밀었더니 선이 너무 돌아 m이 0.28로 줄어듭니다. 6걸음까지 이렇게 천천히 보고, 그 뒤는 빠르게 2만 걸음까지 갑니다.}
\end{frame}

